\begin{figure}[H]
\centering
\begin{tikzpicture}[
    font=\small,
    token/.style={draw=teal!65!black,fill=teal!7,rounded corners=3pt,
        minimum width=2.2cm,minimum height=0.75cm,font=\ttfamily},
    explanation/.style={anchor=west,text width=9.4cm,align=left},
    connection/.style={-{Stealth[length=2mm]},gray!70!black}]
\node[draw=gray!40,fill=gray!6,rounded corners=3pt,
    minimum width=13cm,minimum height=0.8cm] at (5.7,0)
    {\texttt{\detokenize{cl> rfits fintro* "" junk old+}}};

\node[token] (task) at (0.5,-1.2) {rfits};
\node[explanation] (tasktext) at (2.5,-1.2)
    {\textbf{Tarea.} Lee datos de imagen FITS y los convierte
    en imágenes utilizables por IRAF.};
\draw[connection] (task.east) -- (tasktext.west);

\node[token] (source) at (0.5,-2.5) {fintro*};
\node[explanation] (sourcetext) at (2.5,-2.5)
    {\textbf{Entrada: \texttt{fits\_file}.} Selecciona los archivos
    cuyo nombre comienza por \texttt{fintro}; por ejemplo,
    \texttt{fintro0001} y \texttt{fintro0002}.};
\draw[connection] (source.east) -- (sourcetext.west);

\node[token] (selection) at (0.5,-3.8) {\detokenize{""}};
\node[explanation] (selectiontext) at (2.5,-3.8)
    {\textbf{Selección: \texttt{file\_list}.} Para estos archivos
    en disco, la cadena vacía indica leer solo la unidad primaria
    (índice 0), no todas las extensiones.};
\draw[connection] (selection.east) -- (selectiontext.west);

\node[token] (output) at (0.5,-5.1) {junk};
\node[explanation] (outputtext) at (2.5,-5.1)
    {\textbf{Salida: \texttt{iraf\_file}.} Nombre raíz de salida.
    Si se usa esta raíz para varios archivos, se añaden números
    de secuencia. No es un archivo de entrada.};
\draw[connection] (output.east) -- (outputtext.west);

\node[token] (names) at (0.5,-6.4) {old+};
\node[explanation] (namestext) at (2.5,-6.4)
    {\textbf{Nombres originales.} Abrevia \texttt{oldirafname=yes}.
    Intenta usar el nombre guardado en la palabra clave
    \texttt{IRAFNAME} de la cabecera, en lugar de la raíz \texttt{junk}.};
\draw[connection] (names.east) -- (namestext.west);
\end{tikzpicture}
\caption{Desglose del comando \texttt{rfits} para los archivos en disco
 de la práctica. \texttt{cl>} es el indicador del intérprete, no un
 argumento. La recuperación de los nombres originales depende de la
 información disponible en las cabeceras~\cite{iraf}.}
\label{fig:rfits}
\end{figure}
